\scanpage{3}
\mspage{2}
Calculons $\pi_0(Z)$ $(Z=X\times_S Y)$, en utilisant la fibration de Serre
\[
 F\hookrightarrow Z\longrightarrow Y
\]
où $F$ est la fibre de $f:X\to S$, on trouve \inserted{(comme $\pi_0(Y)=1$)} (avec les points base déjà choisis, \uline{et} $z_0=(x_0,y_0)$)
\[
 \pi_0(Z)=\pi_0(F)/\pi_1(Y)\ \text{(opérant à g.)}
\]
d'ailleurs on a, par la fibration de Serre (et $\pi_0(X)=1$)
\[
 F\hookrightarrow X\longrightarrow S
\]
\[
 \deleted{\pi_1(X)\to\pi_1(S)\to\pi_0(F)\to\pi_0\unsure{(X)}\to1}
\]
\deleted{i.e.}
\[
 \pi_0(F)\simeq\pi_1(S)/\underbrace{\pi_1(X)}_{\text{op. à droite}}
\]
d'où, en composant
\[
 \pi_0(Z)=\pi_1(Y)\backslash\pi_1(S)/\pi_1(X)
\]
ce qui montre la bijectivité des $\pi_0$.

\unsure{Récipt.|Reste.} Soit $Z'$ une composante connexe de $Z$. Elle se projette \uline{sur} $Y$, donc $\exists\,z'_0\in Z'_{y_0}$. Son image dans $X$ est \deleted{bien} dans $X_{s_0}=F$, donc
% slanted
\margin{donner une description symétrique en $X,Y$ \ldots}
en la reliant par un arc à $x_0$, on trouve en appliquant $f$ un lacet en $s_0$, soit $\sigma$ [déterminé mod un élément de \inserted{$\mathrm{Im}$} $\pi_1(X,x_0)=G$ par chgt. de l'arc dans $X\ldots$, et modulo un élément de $\mathrm{Im}\,\pi_1(Y,y_0)$ en changeant le point $z'_0$ \ldots.]
% archive sheet number: 2
